\chapter{Frontend, Nutzerrollen und Deployment}

\section{Frontend (Nuxt)}
Die Haupt-UI liegt in \path{frontend/app/app.vue}:
\begin{itemize}
  \item Wallet: Verbinden, Anzeige von Adresse, Chain-ID,
    \emph{VOTE}-Saldo und \emph{Allowance} gegenüber \texttt{PredictionMarket}.
  \item Nicht-Admin: Faucet-Button (\texttt{claimFaucet}), Formular zur
    Markterstellung mit Dauer und Resolution Context---Outcomes sind fest YES/NO.
  \item Marktkarten: Status-Badge, Pools, angezeigte Prozentwerte, Staken mit
    vorherigem Approve für den jeweiligen Betrag.
  \item Admin-Adresse (aus exportierter Konfiguration): Auflösungsbuttons pro
    Markt nach Fristablauf (\texttt{resolveMarket}).
\end{itemize}
Die Composable-Datei \path{useContracts.ts} kapselt alle Lesevorgänge
(\texttt{getMarkets}-Schleife über Märkte und Outcomes) und Transaktionsaufrufe.

\section{Ablauf für Endnutzerinnen}
\begin{enumerate}
  \item Wallet verbinden, Netzwerk lokal oder Sepolia wählen (MetaMask-Einträge
    wie in der README beschrieben).
  \item \emph{VOTE} einmalig fauceten und ggf.\ Saldo aktualisieren.
  \item Markt eröffnen oder bestehenden wählen, \emph{Approve} für gewünschten
    Stake-Betrag, danach YES oder NO staken (solange Zeitfenster und nicht
    aufgelöst).
  \item Nach Auflösung und eigenem Gewinn \texttt{claimReward} auslösen.
\end{enumerate}

\section{Build- und Betriebsszenarien}
\subsection{Lokal (Entwicklung)}
Typische Reihenfolge aus der README:
\texttt{npm run chain} (Terminal~1); \texttt{npm run deploy:local} sowie
\texttt{npm run export:frontend} (Terminal~2); \texttt{cd frontend \&\& npm run dev}
(Terminal~3).

\subsection{Docker \& CI}
Ein Docker-Image bündelt das gebaute Frontend und die exportierten Contract-JSONs,
nicht jedoch Private Keys oder Deployment-Skripte. GitHub Actions
(\path{.github/workflows/}) können das Image erzeugen und z.\,B.\ nach GHCR pushen---
Details im Repo \path{Dockerfile} und \path{docker-entrypoint.sh}.

\section{Schnittstelle zur Chain}
Deployed Adressen und ABIs liegen unter \path{frontend/app/contracts/} (u.\,a.\
\texttt{addresses.json}, Artefakte). Frontend und Dokumentation müssen zur
gleichen Deployment-Generation passen (\texttt{export:frontend} nach jedem Deploy).
